\chapter{Conclusão}

O laboratório atingiu os objetivos propostos. A rede 192.168.50.0/24 foi dividida em
sub-redes /26 e a ausência de comunicação entre sub-redes sem roteador foi demonstrada.
O \textit{firewall} do servidor, com política de negação por padrão e uma única regra
restrita por interface, origem, destino e porta, reduziu a superfície exposta ao cliente
de duas portas (22 e 8080) para apenas a 8080, como comprovado pelas varreduras com
\texttt{nmap}. No armazenamento, o \textit{bucket} configurado propositalmente como público
permitiu leitura, listagem e escrita sem autenticação; após a correção, todas essas
tentativas passaram a ser negadas. O usuário criado consegue apenas listar, ler e gravar
no \textit{bucket} autorizado, e todas as ações fora desse escopo foram negadas.

A principal lição é que controles de rede e de identidade se complementam: cada um cobre
a falha do outro, e a verificação com testes reproduzíveis, com resultado esperado
explícito, é o que permite afirmar que um controle realmente funciona.
